\section{Methodology}
\label{sec:methodology}

Hypothetica is designed as a multi-stage originality assessment pipeline. Given a natural language research idea and a selected evidence source, the method produces an enriched idea representation, retrieves candidate prior work, selects a small set of highly relevant resources, parses their content, and computes an originality score for the submitted idea. The final assessment reports the overall originality score, similarity scores for four originality criteria, and the evidence passages that explain the strongest similarities with prior work.

Originality assessment requires more than surface-level text matching. Two ideas may describe the same technical contribution using different terminology, while a genuinely novel idea may still share domain vocabulary with existing work. Hypothetica therefore combines semantic retrieval, cross-encoder reranking, and large language model reasoning to support concept-level comparison between the submitted idea and prior work. The pipeline is organized into phases with explicit data boundaries so that each step, from retrieval to final score computation, can be analyzed independently.

\subsection{System Overview}

\paperfigureblock[0.88\linewidth][0.36\textheight]{"Hypothetica Flowchart/GeneralOverview.png"}{System Overview Diagram.}{fig:overview}

Hypothetica is structured as a multi-phase pipeline that takes a natural language research idea as input and produces a structured originality assessment as output, as shown in Fig.~\ref{fig:overview}. The pipeline begins with a User Interaction Phase that turns the raw submission into an enriched, well-defined idea description, and continues through a Retrieval Phase that gathers candidate prior art from the user's selected evidence source. A Candidate Selection Phase then narrows this pool down to a small, high-quality shortlist through a combination of vector-based filtering, cross-encoder reranking, and an LLM-based selection step. Once the shortlist is fixed, the Parsing Phase converts each selected resource into structured, searchable text and indexes it in a vector store. The two scoring phases that follow are where the originality verdict is produced: the Candidate Originality Scoring phase analyzes each shortlisted resource independently against the user's idea and produces per-criterion similarity scores together with sentence-level similarity evidence, and the Idea Originality Scoring phase aggregates these per-resource results into a single global originality score, a labeled breakdown across the four originality criteria, and a sentence-level annotation of the user's idea. The final output includes the global originality score, criterion-level similarity values, sentence-level originality labels, and evidence passages that explain the strongest similarity judgments.

The phases communicate through structured data passed from one stage to the next, and the boundaries between them are deliberate: each phase has a single, well-defined responsibility, which makes the pipeline easier to extend with new evidence sources, new scoring criteria, or reporting formats without disturbing the rest of the system. The pipeline exposes parameters for retrieval depth, reranking depth, final resource selection, and final score computation. The default values for these parameters are described in Section~\ref{sec:testing} and summarized in Table~\ref{tab:params}. The remainder of this section describes each phase in detail.

\subsection{User Interaction Phase}

\paperfigureblock[0.48\linewidth][0.50\textheight]{"Hypothetica Flowchart/User-Interaction-Phase.png"}{User Interaction Diagram.}{fig:user_interaction}

The pipeline begins with the User Interaction Phase, which transforms a raw user submission into a structured input that the rest of the system can operate on. As shown in Fig.~\ref{fig:user_interaction}, this phase takes a User Request Payload and produces an Enriched Idea that is passed forward to the Retrieval Phase.

The first step is an Adapter Availability Check. Before any analysis begins, the system verifies that the adapter chosen by the user, such as OpenAlex, Google Patents, or GitHub, is configured and accessible. If the selected adapter is unavailable, the request is rejected. This early validation prevents wasted computation and gives the user immediate feedback when an evidence source cannot be reached.

Once the adapter is confirmed to be available, the Follow-Up Questions step is triggered. A dedicated agent reads the user's research idea and generates three clarifying questions. These questions are designed to cover aspects such as the specific research problem being addressed, the proposed method or approach, and the parts of the idea the user considers novel. The purpose of this step is to surface information that researchers often leave implicit in a short idea description but that is important for targeted retrieval and scoring. Without this clarification, vague ideas tend to lead to less focused retrieval and shallower per-paper analysis later in the pipeline.

The user's responses to these questions are then passed to the Idea Enrichment step. Here, the original idea and the question-answer pairs are combined into a single, structured text block that preserves both the user's initial idea and the additional details provided through the clarifications. This combined text is the Enriched Idea that drives later stages of the pipeline. From this point, retrieval queries, candidate selection, per-paper scoring, and global originality scoring all operate on the enriched description rather than on the raw input.

\subsection{Retrieval Phase}

\paperfigureblock[0.52\linewidth][0.58\textheight]{"Hypothetica Flowchart/Retrieval-Phase.png"}{Retrieval Phase Diagram.}{fig:retrieval}

Once the Enriched Idea Context is produced, the pipeline enters the Retrieval Phase, which is responsible for retrieving a large pool of candidate resources from the selected adapter. As shown in Fig.~\ref{fig:retrieval}, this phase takes the Retrieval Input consisting of the Enriched Idea and the Selected Adapter chosen earlier by the user and outputs N Unique Candidates that are forwarded to the Candidate Selection Phase. The main challenge at this stage is balancing recall and precision. If the search is too narrow, it risks missing important prior work altogether. On the other hand, if it is too broad, it brings back a large number of documents, many of which are not relevant. For this reason, the Retrieval Phase is designed to focus more on recall, while precision is improved later during the filtering steps.

The process starts with Query Variant Generation. Instead of relying on a single query based on the user's idea, a Query Variant Agent produces a structured set of complementary search queries. For general scholarly and patent retrieval, the query set contains a raw query that condenses the core idea, two academic reformulations that separately emphasize the method and the application domain, and a synonym query that captures abbreviations and alternative terms. This structure is intended to support broader recall by covering both natural-language descriptions and the terminology that is likely to appear in titles, abstracts, or patent descriptions.

The query strategy is adapter-aware. For OpenAlex, which performs relevance-ranked full-text search over a broad scholarly index, the agent produces five variants: a raw query, a method-focused academic query, a domain-focused academic query, a synonym query, and a concept-level query that targets related subfields, benchmark names, or survey-level terminology. For GitHub, the agent instead produces four implementation-focused queries built around rare technical terms, framework names, and repository-language expressions, while avoiding generic standalone terms such as ``agent'', ``LLM'', or ``tool''. If query generation fails, the pipeline falls back to a single raw query based on the enriched idea, so retrieval can still proceed.

The Variant Queries List is then sent, one query at a time, to the Adapter Search step. The selected adapter -- OpenAlex, Google Patents, or GitHub -- is invoked through a unified adapter interface that hides the differences between the underlying source services. Each adapter is responsible for issuing the query to its corresponding service, parsing the source-specific response, and converting the raw results into a standardized Resource representation that the rest of the pipeline can operate on. For each query variant, the adapter retrieves candidates up to a configurable retrieval-depth parameter. Operating through this adapter abstraction supports adding new evidence sources without changing the rest of the pipeline.

After all the variant searches finish, the Deduplication step merges their results into a single pool. Since each variant describes the same idea from a slightly different angle, the same resource often shows up in more than one result set, and the system needs to make sure it is only carried forward once. To handle this, deduplication relies on stable identifiers like the OpenAlex Work ID, the patent ID, or the repository's full name, and falls back to DOI when one is available. Anything that survives this step is given a fresh internal ID so the rest of the pipeline can refer to it consistently.

The final output of the Retrieval Phase is a set of N Unique Candidates, where N typically ranges from a few dozen to several hundred depending on the chosen adapter, the retrieval-depth setting, and the number of duplicates removed. N is therefore a run-level output rather than a separate fixed parameter; the default retrieval-depth value is described in the experimental setup and summarized in Table~\ref{tab:params}. This pool is intentionally larger than what would be analyzed in detail, since the goal of this phase is to maximize the chance that any prior work with meaningful similarity to the user's idea is included. Reducing this pool to a small, high-quality set is the responsibility of the next stage, the Candidate Selection Phase.

\subsection{Candidate Selection Phase}

\paperfigureblock[0.52\linewidth][0.58\textheight]{"Hypothetica Flowchart/Candidate-Selection-Phase.png"}{Candidate Selection Diagram.}{fig:candidate_selection}

The Candidate Selection Phase is responsible for tightening this pool down to a small, high-quality set of resources that will be used for detailed analysis. As shown in Fig.~\ref{fig:candidate_selection}, the input to this phase is the Candidate Selection Input, consisting of N Candidates from the Retrieval Phase together with the Enriched Idea, and the output is a fixed number of highly relevant resources that will be passed to the Parsing Phase. The default value is described in the experimental setup and summarized in Table~\ref{tab:params}. The phase is structured as a three-step process that combines vector-based filtering, a more precise but heavier reranking step, and finally a Large Language Model for judgment.

The first stage is Embedding Generation. Each candidate's title and abstract or, in the case of code repositories, the repository description and available documentation are embedded into a dense vector representation using a pretrained text embedding model. Once embeddings are available, the Semantic Similarity Ranking step computes the cosine similarity between the Enriched Idea vector and each candidate vector, and keeps the most similar ones according to a configurable embedding-stage cutoff. This stage removes many low-relevance candidates: candidates that share keywords with the idea but are conceptually unrelated tend to fall away, while candidates that talk about the same underlying problem in different words tend to surface. At this stage the system keeps the candidate set wide rather than aggressive, because the goal is to feed a high-recall input into the more precise reranker.

The second stage is Cross Encoder Reranking. Unlike the embedding step, which encodes the query and each document independently and only compares them at the end, a cross encoder takes the query and a single document together and produces a similarity score in one pass. This usually provides a more precise relevance signal because the model can attend to the interaction between specific words in the query and specific words in the document, but it is also considerably more expensive, which is why it only runs on the top candidates produced by the embedding step. The output of this stage is a top-ranked shortlist whose size is configurable and passed to a language model for the final selection.

The final stage is the Relevance Selection Agent, which acts as an LLM-as-a-judge over the reranked shortlist. Its input is the Enriched Idea and the top-ranked candidates that survived the embedding and cross-encoder stages. For scholarly or patent candidates, the agent receives the title, abstract or description, source metadata, and the available rerank or similarity score; for code repositories, it receives the repository title, description, available documentation content, repository metadata, and the available score. These scores are treated as supporting signals rather than as the final decision rule. The agent is instructed to select the resources that are most useful for a deeper originality analysis by prioritizing conceptual alignment, direct methodological relevance, similar problem framing, and evidence that helps compare the user's proposed approach with prior work. For repository evidence, the criteria are adapted toward implementation quality, feature completeness, code maturity, and direct applicability to the user's idea. The output is a structured list of selected resource identifiers with short selection reasons; if the LLM selection fails or returns no usable identifiers, the pipeline falls back to the top candidates from the cross-encoder ranking.

The output of the Candidate Selection Phase is the final resource list that will be analyzed in detail in the rest of the pipeline. By the time a resource reaches this point, it has passed three filters -- vector similarity, cross-encoder reranking, and language model judgment -- which together provide a stronger relevance signal than any single retrieval step could offer on its own.

\subsection{Parsing Phase}

\paperfigureblock[0.56\linewidth][0.58\textheight]{"Hypothetica Flowchart/Parsing-Phase.png"}{Parsing Diagram.}{fig:parsing}

Before any per-candidate analysis can take place, each of the selected resources needs to be turned into structured, searchable text. This is the responsibility of the Parsing Phase, shown in Fig.~\ref{fig:parsing}. The phase takes the Top-K Candidates produced by the Candidate Selection Phase and ends with an indexed vector store containing the text chunks of every candidate, ready for retrieval during the scoring phases.

The phase begins with Parallel Document Fetching. Each candidate's available full content is collected concurrently from its source. PDFs are retrieved for OpenAlex, the Google Patents API is queried directly for patent descriptions, and the README content already collected during retrieval is reused for GitHub repositories. The fetched content is then converted into a Markdown representation so that the rest of the pipeline can treat resources from different adapters in the same way.

Once a uniform document is available, the Extract Section Content step parses out the structural skeleton of each resource. Headings are identified and the text under each heading is associated with that heading as its section content. Sections that contribute little to originality assessment, such as references, acknowledgments, and patent-specific legal sections, are skipped at this stage so that they do not pollute the downstream representation.

The remaining steps are repeated for each candidate. The Document Chunking Processor splits each section into smaller chunks, with target sizes chosen to fit the context windows of both the embedding model and the language model. The Vector Store Ingestion step then embeds these chunks and stores them with their source metadata. The output of the Parsing Phase is an indexed vector store ready to support retrieval during the Candidate Originality Scoring and Idea Originality Scoring phases.

\subsection{Candidate Originality Scoring}

\paperfigureblock[0.72\linewidth][0.54\textheight]{"Hypothetica Flowchart/Candidate-Originality-Scoring-Phase-Layer-1.png"}{Candidate Originality Scoring Diagram.}{fig:candidate_scoring}

The Candidate Originality Scoring phase is where each candidate is independently compared against the user's idea to produce a structured similarity profile. As shown in Fig.~\ref{fig:candidate_scoring}, the input to this phase is the Enriched Idea together with the processed candidates and the indexed candidate text, and the output is a per-candidate Candidate Originality Scoring Result containing a similarity score, criteria-level scores, sentence-level analysis, matched sections, and a justification for each score. The phase loops over every candidate and treats each one as a separate scoring task, so that a highly similar resource remains visible in the later aggregation stage.

For each candidate, the system first builds a focused textual representation of the resource. The Extracted Fixed Sections step pulls out parts of the document that are generally useful for originality assessment, such as the abstract, introduction, and conclusion when they are available. In parallel, the Contextual RAG Retrieval step queries the vector store with the Enriched Idea and retrieves the chunks from the same candidate that are most semantically aligned with it. Combining a fixed structural skeleton with retrieval-driven context lets the language model see both the standard parts of the resource and the parts that specifically discuss similar problems or methods, without overwhelming the prompt with the entire document.

This combined context is then passed to two parallel evaluation tracks. The first track is the Multi-Criterion LLM Evaluation, which scores the candidate against the Enriched Idea on four originality dimensions: Problem Similarity, Method Similarity, Domain Similarity, and Contribution Similarity. Problem Similarity measures whether the candidate addresses the same research question or gap. Method Similarity measures whether it uses the same technical approach, architecture, or pipeline. Domain Similarity measures whether it operates in the same application area, task setting, data type, or target system. Contribution Similarity measures whether the candidate makes the same substantive claim or demonstrates the same proposed advance. Each criterion is scored in a separate LLM call with its own rubric and 5-point Likert definitions. This keeps the comparison focused: the model evaluates one dimension at a time instead of balancing several partly related judgments in a single response. Each call also returns a short justification tied to specific candidate content, which is later used as evidence in the report.

The second track is the Sentence-Level Similarity Analysis, which examines the user's idea at a finer granularity. Instead of evaluating the idea as a whole, this step scores every individual sentence of the Enriched Idea against the candidate, again using the same four criteria and the same 5-point Likert scale. For each sentence, the LLM returns criterion-specific matched sections from the candidate, including the section heading, a short evidence snippet, the criterion it supports, its similarity level, and a brief reason. These sentence-level results are then checked against the candidate-level scores. A sentence-level criterion score that falls more than one step below the candidate-level score for the same criterion on the five-point scale is set to zero, and its matched sections are removed. This keeps the sentence evidence tied to the broader candidate judgment instead of allowing isolated sentence matches to dominate the analysis.

The Likert scores produced by the two tracks are then mapped to floats in the $[0, 1]$ range, where 1 corresponds to near-complete similarity and 0 corresponds to no meaningful similarity. The Similarity Score Calculation step combines the four candidate-level criteria scores into a single per-candidate similarity score using a weighted average:

\begin{equation}
	\text{sim}_{\text{candidate}} =
	w_p \cdot S_{\text{problem}}
	+ w_m \cdot S_{\text{method}}
	+ w_d \cdot S_{\text{domain}}
	+ w_c \cdot S_{\text{contribution}}
\label{eq:sim_candidate}
\end{equation}

Here, $S_{\text{problem}}$, $S_{\text{method}}$, $S_{\text{domain}}$, and $S_{\text{contribution}}$ are the normalized similarity scores for the four originality dimensions. The weights $w_p$, $w_m$, $w_d$, and $w_c$ control how strongly each dimension contributes to the candidate-level similarity score. Contribution similarity is treated as the primary novelty signal because it captures whether the candidate already makes the same substantive claim; method similarity is the second strongest signal because it reflects whether the same technical approach has been applied. Problem and domain similarity provide supporting context but carry lower weight, since the retrieval step already constrains candidates to the same research area. The default weight values used in the experiments are given in the experimental setup and summarized in Table~\ref{tab:params}.

The output of this phase, for each candidate, is a Candidate Originality Scoring Result containing the per-candidate similarity score, the four criteria scores together with their justifications, the sentence-level similarity analysis, and the matched evidence sections retrieved from the candidate. These results are the structured input to the next phase, where they are aggregated into a single global originality score for the user's idea.

\subsection{Idea Originality Scoring}

\paperfigureblock[0.86\linewidth][0.42\textheight]{"Hypothetica Flowchart/Idea-Originality-Scoring-Phase-Layer-2.png"}{Idea Originality Scoring Chart.}{fig:idea_scoring}

The Candidate Originality Scoring phase produces an independent similarity profile for every Top-K candidate. The Idea Originality Scoring phase takes this collection of per-candidate results and turns it into a single structured assessment of how original the user's idea is. As shown in Fig.~\ref{fig:idea_scoring}, the input to this phase is the set of scoring results for every candidate, and the output is the Idea Originality Scoring Result used in the final report. The phase has three parallel branches -- Aggregated Criteria, Global Similarity Computation, and Sentence Annotations -- followed by a final originality score derived from the global similarity.

The Global Similarity Computation branch determines the final originality assessment. Each candidate already carries a per-candidate similarity score $\text{sim}_i$ produced in the previous phase, where $i$ indexes the $N$ selected candidates. A simple average of these scores would dilute a strong similarity signal whenever the retrieval pipeline returned a few weakly related candidates alongside a genuinely close match. To preserve that signal while still being robust to noise, the global similarity is computed as a blend of the maximum and the average of the per-candidate scores:

\begin{equation}
\begin{alignedat}{2}
	\text{sim}_{\text{global}}
	&= w_{\max} \cdot \max_{i}(\text{sim}_{i})
	&\quad{}+{}& \frac{w_{\text{mean}}}{N} \sum_{i=1}^{N} \text{sim}_{i}
\end{alignedat}
\label{eq:sim_global}
\end{equation}

The two weights $w_{\max}$ and $w_{\text{mean}}$ sum to one and are configurable. The maximum term gives priority to the closest prior work, while the mean term keeps the score sensitive to the broader candidate set. As a result, the global similarity depends mainly on the strongest match without ignoring the remaining candidates.

The global similarity is then converted to a final originality score on a 0--100 scale:

\begin{equation}
	\text{originality} = \left(1 - \text{sim}_{\text{global}}^{p}\right) \cdot 100
	\label{eq:originality}
\end{equation}

The exponent $p$ controls how quickly originality decreases as similarity increases. A linear conversion would reduce the originality score at a constant rate, but this formulation treats the high-similarity end of the scale as more sensitive: once a candidate is close to the submitted idea, small additional increases in similarity have a larger effect on the final score. Using $p > 1$ captures this behavior while still preserving the ordering of candidates by similarity. The resulting integer score is then mapped to one of three labels -- high, medium, or low originality -- using configurable thresholds.

In parallel with the global computation, the Criteria Breakdown branch produces a single reported value per criterion across all candidates. For each criterion separately, letting $c_i$ denote the score of candidate $i$ on that criterion, the system combines the per-candidate scores using the same blend of maximum and mean as the global similarity:

\begin{equation}
\begin{alignedat}{2}
	S_{\text{breakdown}}
	&= w_{\max} \cdot \max_{i}(c_i)
	&\quad{}+{}& \frac{1 - w_{\max}}{N} \cdot \sum_{i=1}^{N} c_i
\end{alignedat}
\label{eq:breakdown}
\end{equation}

For each criterion, the reported value therefore reflects both the strongest candidate-level similarity and the average similarity across the analyzed candidates. Using the same $w_{\max}$ weight as the global similarity keeps the criterion summary consistent with the blend that drives the final score. These breakdown scores are not used in the final originality score; they indicate which dimensions of the idea are most similar to prior work.

The Sentence Annotations branch carries these evidence links into the final result. For every sentence in the Enriched Idea, the system collects the retained sentence-level criterion scores from all analyzed candidates. For each criterion separately, it keeps the highest score observed across candidates. When a new maximum is selected, the system also stores the strongest matched section for that criterion from the same candidate; if no criterion-specific passage is available, it falls back to the highest-scoring matched section for that sentence. The retained criterion scores are converted back into Likert levels and mapped to high, medium, or low originality labels using configurable thresholds. The overall sentence label is the least original label assigned to any of its criteria, and the sentence is linked to the evidence passages that produced those criterion-level labels. These sentence annotations explain the final report but do not affect the global originality score.

The three branches are then combined into the Idea Originality Scoring Result, which contains the final originality score and label, the aggregated criteria values, the per-sentence annotations with their linked evidence, and a short natural-language summary generated by a dedicated Summary Agent. The summary takes the global score, the aggregated criteria, and the distribution of sentence labels as input, and produces a one-to-two-sentence explanation that highlights the main areas of similarity and offers an actionable recommendation. This final structured object forms the originality report returned by the pipeline. The following section describes how these components are realised in the system's implementation.
